% !TEX root = muni_tax_sorting.tex
% !TEX program = pdflatex
% !BIB program = biber
\documentclass[12pt]{article}
\usepackage[margin=1in]{geometry}
\usepackage{setspace}
\doublespacing
\usepackage{graphicx}
\usepackage{booktabs}
\usepackage{amsmath}
\usepackage{amssymb}
\usepackage{caption}
\usepackage{float}
\usepackage{hyperref}
\usepackage[style=apa,backend=biber,natbib=true]{biblatex}
\addbibresource{local_income_taxes_2.bib}
\usepackage{bm}

\newcommand{\sym}[1]{\textsuperscript{#1}}

\hypersetup{
    colorlinks=true,
    linkcolor=black,
    citecolor=black,
    urlcolor=blue
}

\title{Do Local Taxes Cause Residential Sorting? \\
\large Income versus Property Taxes in Pennsylvania}
\author{Matthew Hockert\thanks{University of Minnesota. Email: hocke069@umn.edu}}
\date{\today}

\begin{document}

\maketitle

\begin{abstract}
Do local taxes cause residential sorting, and does the answer depend
on which tax? I compare two local taxes that fall on the same
Pennsylvania municipalities over 2002--2019: the Earned Income Tax
(EIT), levied on the mobile worker, and the property tax, levied on
the immobile dwelling. Both are estimated at the municipality level in
a continuous stacked difference-in-differences design using block-level
LEHD/LODES data aggregated to municipalities. The EIT follows a
piecewise max-rule, which lets me split otherwise-identical rate
increases into a binding cell that moves the worker's bill and a
within-design placebo cell that does not; a typical binding increase
reduces the count of workers who both live and work in their own
municipality by roughly $13\%$ ($-0.573$ log points per percentage
point, $\approx -43\%$), with a clean realized-dose pre-trend, while the
placebo cell is null. To ensure the income-tax response is not merely proxying for broader
local fiscal burden, I include the municipality's property-tax burden
among the controls; the binding-cell estimate is unchanged by its
inclusion. The mechanism is intuitive: the EIT follows the worker to the
municipality of residence, so relocating avoids it, whereas a property
tax attaches to the parcel and cannot be escaped by moving within the
same labor market.
\end{abstract}

\newpage

%==========================================================
\section{Introduction}
\label{sec:intro}
%==========================================================

The Tiebout hypothesis predicts that residents sort across
jurisdictions which offer different fiscal amenities at different costs and residents will 'vote with their feet' and choose specific bundles of amenities for a given amount of taxation, this is called the Tiebout margin \citep{tiebout1956}. A survey of the literature conducted by \citet{agrawalLessonsStateLocal} highlights the current findings of state and local income taxation and identifies sub-state (county, municipality, school districts) income-tax sorting as one of the field's open questions.
The literature has focused mostly on state tax adjustments on high-earners
\citep{young2016, agrawalRelocationRichMigration2019}. The research shows that
top-income households are generally mobile when rates between states are large, and \citet{agrawal2015gradient} documents sharp discontinuities among commodity-taxes at sub-state borders. The closest existing research comes from
\citet{yangRESPONSIVENESSINCOMELOCAL2017} which  finds small migration effects for county-level income tax adjustments in Indiana.

This paper estimates municipal-level migration effects among workers in Pennsylvania from 2002--2019, by merging local income tax data with the LEHD Origin-Destination
Employment Statistics (LODES), which links each worker's residence to their workplace at annually and at sub-municipal geography. In Pennsylvania there are roughly $2{,}500$ municipalities and $500$ school districts (SD) that set
their own Earned Income Tax (EIT) rates
\citep{bagchiEffectsIntroducingWithholding2022,agrawalLessonsStateLocal},
producing a deep cross-section of multiple sub-state income-tax opportunities. The EIT offers two additional dimensions of variation that my
empirical design exploits. First, the EIT is set independently by two
overlapping rate-setting jurisdictions (municipalities and school
districts) and a resident's home-side tax bill is the sum of the
municipal resident rate and the school-district rate. This paper's primary
focus is the municipal resident rate while the school district rate enters the empirical
design but is treated as secondary.
Second, the EIT operates using a piecewise maximization rule where each worker's bill is
$T = \max(\tau^{\text{home}}, \tau^{\text{work,muni}})$, where
$\tau^{\text{home}} = \tau^{\text{res,muni}} + \tau^{\text{SD}}$ is the
sum of the municipal resident rate and the school-district resident
rate at the worker's address (the home stack), and $\tau^{\text{work,muni}}$ is the
work rate of the worker's workplace municipality. The workplace rate is also referred to as the nonresident rate, however, I refer to it as the work rate because the primary outcome variable of is the number of people that live and work within a municipality, which is the only variable that has the cleanest causal identification to the work and home-stack rate changes A statutory home-side rate
change only moves the bill if the home stack is on the binding side
of the max. Act~32 of 2008 (effective 2012) consolidated EIT
collection into county-based Tax Collection Districts and mandated
employer withholding at the higher of the two rates, which made the
max-rule mechanically enforced.
\citet{bagchiEffectsIntroducingWithholding2022} estimates that Act~32 resulted in a $14 \%$ increase in compliance and much higher local revenues. My paper does not analyze compliance but rather estimates
residential-sorting and provides additional analysis pre-post Act~32.

Beyond contributing to the general literature on local tax rate changes and residential sorting  this paper's methodological contribution is utilizing the piecewise structure of the EIT to generate a binding decomposition to test how rate changes that move the tax bill from the work side to the home side effect residential sorting. A statutory rate change moves the worker's bill only when the home stack is on the binding side of the
maximization function, but that same statutory change produces a predicted-null when
the home stack stays below the work rate. Splitting
rate-change events by whether or not they are binding isolates
whose bill actually moves and those whose bill does not. This within-design heterogeneity is what \citet{agrawalLessonsStateLocal} identifies as missing from the existing U.S. local-income-tax literature.
The decomposition is implemented at the municipality-year level in
a continuous stacked difference-in-differences design that uses
never-changed municipalities as the control pool. The
school-district rate enters the design as a co-treatment on the
home-side stack but is not used as a separate unit of analysis.
Section~\ref{sec:theory} formalizes the binding decomposition and
lays out the worker-bill differencing function each cell implies.

The design shows that when the municipality-level tax is increased on an already binding home-stack (Inc-bind) residents who live and work there tend to decline. At the average Inc-bind dose of $0.24$ pp cumulative
over five years, the muni-rate change reduces the log count of
workers who both live and work in the municipality by roughly
$13\%$ ($\beta = -0.573$ per pp, $0.24 \times -0.573 = -0.14$ log
points). The increase becomes (Inc-becomes) cell, where an increase pushes a
previously non-binding home stack above the work rate produces a
smaller per-pp estimate of $-0.250$, implying cumulative dose of $0.41 \times -0.250 =
-0.10$ log points (about $-10\%$). The Inc-nobind cell is a null
within-design placebo, indicating that the mechanism is producing estimates that match the hypothesized effects.

The second contribution is comparative. The same municipalities levy a
property tax---municipal and county real-estate millage---alongside the
EIT, and I estimate its sorting effect in the same municipality-level
design. The two taxes differ in what they fall on: the EIT is levied on
the mobile worker, the property tax on the immobile dwelling, so the
classic capitalization argument \citep{oates1969} predicts a much
smaller mobility response to the property tax, with the burden absorbed
into house prices instead. To make the treatment a policy object rather
than a mechanical one, I use the effective property-tax rate---statutory
millage times the county common-level ratio of assessed to market
value---and isolate its policy-driven (discretionary millage) component,
which strips out the assessment drift that would otherwise confound the
treatment with local house-price dynamics. Each leg controls for the
concurrent change in the other tax, so the income- and property-tax
elasticities are separately identified.

Because the property tax is identified far less cleanly than the EIT---it
changes nearly every year, so there is no never-treated control group and
no non-binding cell, leaving the realized-dose pre-trend as its only
check---I treat it as a control on the income-tax response rather than as
a second treatment estimated in its own right. Including the
municipality's property-tax burden alongside the EIT leaves the
binding-cell estimate unchanged, so the sorting I measure is a response to
the income tax specifically, not to local fiscal burden in general.

The remainder of the paper is organized as follows.
Section~\ref{sec:background} describes the institutional setting---the
EIT and Act~32, and the municipal and county property tax.
Section~\ref{sec:theory} introduces the six-cell binding decomposition
for the EIT. Section~\ref{sec:data} describes the data sources: LODES,
the DCED rate and millage registers, the STEB common-level ratios, and
the ACS income and home-value series. Section~\ref{sec:strategy}
explains the continuous stacked difference-in-differences design at the
municipality level. Section~\ref{sec:results} reports the
EIT binding decomposition and the EIT revenue response.
Section~\ref{sec:robustness} reports robustness diagnostics.

%==========================================================
\section{Institutional Background}
\label{sec:background}
%==========================================================

Pennsylvania has a deep cross-section of jurisdictions setting
their own income tax rates with approximately 2{,}500 municipalities and 500 school
districts each levy a separate Earned Income Tax on essentially
all earned income earned in or by Pennsylvania residents
\citep{bagchiEffectsIntroducingWithholding2022,agrawalLessonsStateLocal}.
Unlike property taxes, which are paid by owners and capitalized into
home values, the EIT falls directly on workers and varies across
municipal boundaries, creating
sharp spatial discontinuities in tax burden that can influence where
workers choose to live and work.

\subsection{Pennsylvania Earned Income Tax}
\label{subsec:pennsylvania}



Pennsylvania's EIT is mandatory for all municipalities under the
Local Tax Enabling Act of 1965 \citep[Act 511;][]{paAct511_1965} and applies to all earned income. The standard combined cap is 1\%, typically split as
0.5\% municipal and 0.5\% school district. Home Rule municipalities
may exceed this cap \citep{paHomeRule}, and Philadelphia operates
separately under the Sterling Act of 1932
\citep{paSterlingAct_1932} at approximately 3.9\%. EIT is a valuable revenue funder for local governments next to property-taxes. However, quantifying the effect of property tax millage on sorting would
require a 2002--2019 municipality-level millage panel that is not currently
available from DCED, and I leave the comparison to property-tax
sorting as a separate exercise. Out of caution that property taxes may impact sorting I control for the share of revenue that a municipality collects from property taxes, EIT, and transfers within a given year.



Each Pennsylvania municipality has three EIT rates relevant to a
worker:
%
the municipality resident rate $\tau^{\text{res,muni}}$ levied on residents of the municipality, the municipality work rate $\tau^{\text{work,muni}}$ levied on workers employed in the municipality, and the school-district rate $\tau^{\text{SD}}$ levied on residents only (school districts cannot tax nonresidents).
%
Each address is associated with a six-digit Political Subdivision (PSD)
code identifying its (municipality, school district) pair. Workers report this
code on the Residency Certification Form, and employers use it to
determine withholding obligations.

%==========================================================
\begin{figure}[H]
  \centering
  \caption{Timeline of Pennsylvania local-tax statutes; the shaded band is the 2002--2019 estimation window.}
  \label{fig:law_timeline}
  \includegraphics[width=0.95\textwidth]{../figures/law_timeline.png}
\end{figure}
%==========================================================

Under Act 32 (effective January 1, 2012), a worker's total EIT bill
is determined by the maximum of the home-side rate stack and the
work-location rate, and the bill is split piecewise
between home and work jurisdictions:
%
\begin{equation}
\bigl(\tau^{\text{home}},\ \tau^{\text{work}}\bigr) =
\begin{cases}
\bigl(\tau^{\text{res,muni}}_j + \tau^{\text{SD}}_j - \tau^{\text{work,muni}}_i,\ \tau^{\text{work,muni}}_i\bigr)
   & \text{if } \tau^{\text{res,muni}}_j + \tau^{\text{SD}}_j > \tau^{\text{work,muni}}_i \\[4pt]
\bigl(0,\ \tau^{\text{work,muni}}_i\bigr)
   & \text{if } \tau^{\text{res,muni}}_j + \tau^{\text{SD}}_j \leq \tau^{\text{work,muni}}_i
\end{cases}
\label{eq:piecewise}
\end{equation}
%
where $j$ is the worker's home municipality and $i$ is the work municipality. The
work jurisdiction always collects its full work rate first, then
the home municipality and school district split whatever portion of the
worker's bill is left over after the work claim. If the work
rate exceeds the home stack, the home jurisdictions get nothing. The
worker's total bill is the higher of the two stacks:
%
\begin{equation}
\text{Bill}_{ji} = \max\!\bigl(\tau^{\text{res,muni}}_j + \tau^{\text{SD}}_j,\ \tau^{\text{work,muni}}_i\bigr).
\label{eq:bill}
\end{equation}

Equation~(\ref{eq:piecewise}) implies which margin a worker can use to
respond to a rate increase. When the home stack binds
($\tau^{\text{home}}_j > \tau^{\text{work,muni}}_i$), the worker can lower
the bill only by moving their residence to a lower-rate municipality or school district.
When the work rate binds, the worker can lower the bill
only by changing jobs to a lower-rate work location, while moving homes
without changing jobs has no effect. For workers who live and work in
the same municipality, the home stack and work rate are both set by that municipality, and the bill is determined by the home stack on the binding side.

A worker who lives in municipality $j$ and moves to a lower-rate municipality $j'$
while keeping the same job at $i$ saves at least the SD rate at $j$:
the new bill is $\max(\tau^{\text{home}}_{j'}, \tau^{\text{work,muni}}_i)$,
which is at most the old bill $\tau^{\text{home}}_j$ minus
$\tau^{\text{SD}}_j$. The empirical design follows this logic directly by using both the treatment
$\Delta\tau^{\text{muni}}$ and the binding classification 
built from the worker's home stack. The regressions therefore measure
sorting in response to the same tax a mover would actually face.
This paper does not estimate the effect of nonresident (workplace tax) increases on sorting due to the lack of binding rate changes of the workplace tax over the home-stack. For the work side to bind, a single rate-setter must exceed the home-stack, which under the standard Act 511 rate structure almost never occurs: across 2002–2019, the work rate is the binding side in roughly $0.2\%$ of municipality-years. The handful of exceptions are municipalities with statutory authority to exceed the Act 511 cap (Home Rule and Act 47 distressed cities). Because there are few work-binding events, a work-side binding decomposition is not identifiable.

Philadelphia is the one exception to the Act 32 framework. It operates under
the Sterling Act of 1932 where Philadelphia residents pay the city resident
rate regardless of work location with no offset and Philadelphia
nonresidents pay the city work rate, with no residual flowing
to home jurisdictions in other Pennsylvania municipalities. My main estimates
retain Philadelphia in the sample for completeness and the results are robust
to dropping it.

Act 32 transformed the system by consolidating collection into
approximately 67 county-based Tax Collection Districts (TCDs), each with a single appointed collector.
The law mandated standardized employer withholding at the higher of
the employee's resident and work rate. This had two consequences.
First, it made the EIT mechanically binding where the rate changes now
appear in every worker's paycheck through employer withholding,
regardless of whether the worker resides in the taxing municipality.
Second, it eliminated the possibility of avoidance
of the work-location claim through weak enforcement.
After Act 32, a worker could no longer effectively avoid the
work-rate portion owed to the work municipality.

However, Act 32 did not eliminate all tax savings from
residential relocation. The piecewise rule in
Equation~(\ref{eq:piecewise}) continues to imply that a worker whose
original home stack binds can reduce the bill by moving to a
sufficiently lower-stack home jurisdiction where the new bill is
$\max(\tau^{\text{home}}_{j'}, \tau^{\text{work,muni}}_i)$, which falls
below the old home stack $\tau^{\text{home}}_j$ whenever
$\tau^{\text{home}}_{j'} < \tau^{\text{home}}_j$ and the original
home stack was the binding side. Act 32 thus changed enforcement
and made the max-rule mechanically binding through withholding; it
did not eliminate the residential-relocation margin that the
mechanism in this paper exploits.
\citet{bagchiEffectsIntroducingWithholding2022} highlighted the impact of Act 32 by showing that consolidated employer
withholding increased PA EIT compliance thus raising reported earned
income in the rate-raising municipality. The present paper complements that work by isolating
the residential-sorting response to the same withholding-enforcement
mechanism exploiting the variation in whether a given
rate change actually moves the worker's bill (the binding
decomposition); the pre/post-Act-32 split in
Section~\ref{sec:robustness} shows the binding-cell sorting
response is strengthened (the effect existed prior to Act 32) after Act 32, complementing the
compliance-response finding.

Figure~\ref{fig:act32_cdf} documents the institutional break at
Act 32. Before 2012, fewer than 60\% of PA municipalities reported a
work rate to DCED; after 2012, when consolidated withholding
made the max-rule mechanically binding, coverage rises sharply.
This is the pre-/post-Act-32 enforcement break used in the
heterogeneity test of Section~\ref{subsec:rob-act32}.

\begin{figure}[H]
    \centering
    \includegraphics[width=0.85\textwidth]{../figures/act32_cdf.png}
    \caption{Share of PA municipalities reporting a work
    rate to DCED, by year. The break at 2012 reflects the Act 32
    requirement that employers withhold at the higher of the home
    resident and work rates.}
    \label{fig:act32_cdf}
\end{figure}

School districts in Pennsylvania also levy an EIT under Act 511
\citep{paAct511_1965}, typically at 0.5\% (the other half of the
standard 1\% cap). Act 1 of 2006 \citep{paAct1_2006} (the Taxpayer
Relief Act) additionally authorized
school districts to exceed the 0.5\% cap specifically to fund
property tax relief. I find that the correlation between the municipal and school-district resident rates is only weakly negative of $-0.12$ indicating the two taxing jurisdictions move on relatively separate schedules. Since every census block sits in both a municipality
and a school district, workers face two overlapping EIT rates
that can change independently, providing a natural decomposition
of the residential sorting response into municipal and school
district components.

\subsection{Pennsylvania Property Tax}
\label{subsec:proptax-inst}

The same municipalities, together with their counties, levy a real-estate
property tax. Each taxing body sets a millage rate applied to the
assessed value of a parcel; the assessed value is a county-set
fraction of market value, and the State Tax Equalization Board (STEB)
certifies that fraction---the common-level ratio (CLR)---for each county
annually. The municipal and county millage are reported in the same DCED
tax register that reports the EIT, so a municipality-by-year panel of
both is directly available; the school-district real-estate millage is
also reported but the school district is not used as an estimating unit
here. Two institutional features shape the analysis. Pennsylvania
reassesses infrequently and irregularly---many counties had not reassessed
in decades during the sample, and only $24$ of $67$ reassessed at all over
2002--2019---so the CLR drifts downward between reassessments as market
values outpace stale assessments. And the Act~1 anti-windfall provisions
require that millage be rolled back to revenue-neutrality at a
reassessment, with any real increase requiring a separate, capped vote.
These features mean the statutory millage and the effective burden it
implies can diverge, which is why the analysis works with the effective
rate (millage $\times$ CLR) and its policy-driven component rather than
raw millage. Unlike the EIT, the property tax falls on the dwelling
rather than the worker.



\section{Data}
\label{sec:data}
%==========================================================

\subsection{Data Sources}
\label{subsec:sources}
\label{subsec:sample}
\label{subsec:pipeline}

The analysis merges the LEHD  residence--workplace dataset and the  Pennsylvania EIT rate records to create a panel of rate changes and workforce data over space and time. The two datasets do not share common identifiers, and was bridged using TIGER~2020 block geometry, which
nests blocks within county subdivision code (cousubs). A
municipality-name crosswalk maps each DCED municipality ID to its
corresponding COUSUBFP code, matching $2{,}496$ of $2{,}505$ PA municipalities
($99.6\%$). The remaining $9$ are small boroughs that have either dissolved or
merged during 2002--2019 and are dropped from the analysis.

The residence--workplace data come from LEHD LODES Version~8,
2002--2019 \citep{lehdLODES}. The Census Bureau applies
synthetic noise infusion at the block level to protect respondent
confidentiality and is calibrated so that aggregation to tract and larger
geographies approximates ground truth
\citep{grahamDESIGNCOMPARISONLODES}. The sample period of 2002--2019
reflects LODES coverage beginning in 2002 and
2020--2023 to not include disruptions in the data resulting from the COVID Pandemic.

Municipal and school-district EIT rates come from the Pennsylvania
Department of Community and Economic Development (DCED) MunStats
portal \citep{padcedMunStats}, scraped annually. In some DCED reports, the flat-dollar Local Services Tax is filled in the same field
as the EIT in some years, so rates above $5\%$ are treated as
data-entry errors and dropped except for Home Rule cities and Philadelphia.

Municipal fiscal controls come from the PA DCED Statewide Municipal
Annual Financial Report database \citep{padcedAFR}. School-district
fiscal controls come from the Census Annual Survey of School System
Finances (F-33) \citep{censusF33}; the NCES \texttt{LEAID}-to-DCED
SD ID crosswalk covers all $500$ school districts. All fiscal
controls are share-based and time-varying. Municipality regressions
use six DCED AFR shares: three expenditure shares
(police-services, streets, debt-service) and three revenue shares
(EIT-revenue, real-estate-tax, intergovernmental-revenue).
LODES blocks are aggregated to muni for analysis.

The municipality--SD assignment uses the modal SD ID per muni in
the DCED rate panel ($97.6\%$ of municipalities are single-SD), so
each muni-year carries a single $\tau^{\text{SD}}$ co-treatment.
Municipality regressions
further impose a
minimum employment threshold of $50$ total jobs, which protects the
share-based outcomes from mechanical noise (a single worker moving
changes a share by $1/N$, so at $N=10$ a one-worker move dominates
the per-pp coefficient) and matches the LODES synthetic-noise
calibration that is validated at tract scale and above
\citep{grahamDESIGNCOMPARISONLODES}. The threshold drops $185$ of
$2{,}559$ municipalities ($7.2\%$), almost all sparsely populated
NC-pool townships; the treated cohort is essentially unaffected.

For the property tax, municipal and county real-estate millage come
from the same DCED MunStats register, scraped annually
\citep{padcedMunStats}; the county tax register reports general-purpose
and total millage for all $66$ assessing counties (Philadelphia is a
consolidated city--county). County common-level ratios come from the
Pennsylvania Department of Revenue's historical CLR factor series, in
which the published factor is the reciprocal of the CLR; its
assessment-base-change footnotes also date each countywide reassessment,
which I use to flag reassessment years. The effective property-tax rate
combines municipal and county millage with the county CLR; across
municipality-years its median is about $2.1\%$ of market value (with the
school district excluded as an estimating unit).

\subsection{Descriptive Statistics}
\label{subsec:descriptive}

This section presents the summary statistics at the municipality
level. The primary outcome variables are resident-worker counts,
workplace counts, in-commuter flows, and the live-and-work share.
Fiscal outcomes used as control variables (log EIT revenue, log
total tax revenue, log total revenue, log police expenditure) are
reported in Appendix~\ref{app:robustness}.

Table~\ref{tab:summary_muni} summarizes the estimation panel.

% (block summary removed; muni is the only estimating unit)
% \input{../tables/summary_block.tex}

\input{../tables/summary_muni.tex}

Table \ref{tab:summary_muni} presents the municipality level summary statistics. Table \ref{tab:summary_muni} shows that there are roughly $43{,}374$ muni-year observations, and the resident rate is tightly clustered
near the Act~511 default at $\bar{\tau}^{M}=0.53\%$ and a standard deviation of $0.16$pp. There is a long right tail indicated by the max rate of $4.33\%$ concentrated in Philadelphia and
a handful of Home-Rule and Act~47 distressed cities. The
work rate has $34{,}080$ observation, about $9{,}300$ fewer
observations than the resident rate, due to DCED's lack of coverage for the work-rate field being incomplete before Act~32 mandated
reporting. The SD rate has nearly twice the
muni-rate dispersion (SD $0.29$pp vs.\ $0.16$pp) and maxes at $2.05\%$, potentially making the school-district
rate-setter a larger statistical lever. The home stack averages
$1.15\%$, above the $1\%$ Act~511 cap, because Act~1 SD increases
and Home-Rule muni rates push the upper level past the ceiling.
Only $9.5\%$ of muni jobs are
held by residents of the same muni showing that residential sorting operates
on a thin within-muni margin.



% \input{../tables/summary_sd.tex}  % (SD summary table removed; see paper scope)

\input{../tables/cell_counts_summary.tex}

Table~\ref{tab:cell_counts_summary} reports the cell
assignments on the muni resident-rate side. Of roughly $2{,}560$
municipalities, after dropping those that decrease or flip-flop any of
their three rates, $1{,}687$ never-changers form the NC control pool and
only $77$ enter one of the three increase cells, split as $24$ Inc-bind,
$39$ Inc-becomes, and $14$ Inc-nobind.
Inc-becomes shows a larger average
dose than Inc-bind, indicating that crossing the binding threshold typically
requires a bigger statutory move than slowly inching up an already-binding
rate.

Figure~\ref{fig:rate_changes_by_year} reports the timing of
municipal resident-rate changes across the panel. Increases jumped in the early 2000s, from 2002 to 2009 and taper off after. Most of the identifying variation in my models predates Act 32, highlighting the significant effects found are on smaller increases and further warranting a pre-post 2012 decomposition.

\begin{figure}[H]
    \centering
    \includegraphics[width=0.95\textwidth]{../figures/rate_changes_by_year.png}
    \caption{Number of PA municipalities changing their resident
    rate, by year. Increases shown above zero; decreases shown below.
    Dashed line: 2012, when Act 32 consolidated withholding.}
    \label{fig:rate_changes_by_year}
\end{figure}

Lastly, figure~\ref{fig:pa_map} shows the spatial distribution of the rate changers in my sample. Most rate changers occur along the east coast, but some are also scattered within the interior of the state. This spatial distribution highlights the value of future work to understand the effects when two neighboring municipalities raise or lower rates across time.

% --- Treatment Geography Map ---
\begin{figure}[H]
    \centering
    \includegraphics[width=0.85\textwidth]{../figures/pa_map.png}
    \caption{Pennsylvania EIT rate changes, 2000--2019. Red =
    increased; blue = decreased; purple = both; grey = no change.}
    \label{fig:pa_map}
\end{figure}

% ---- Design schematics parked here for assessment (2026-07-07) ----
\begin{figure}[H]
  \centering
  \caption{The max-rule: when does a rate increase move the bill?}
  \label{fig:binding_schematic}
  \includegraphics[width=0.95\textwidth]{../figures/schematic_binding_mechanism.png}
\end{figure}

\begin{figure}[H]
  \centering
  \caption{Stacked construction: one sub-experiment per cohort; never-changers replicated as controls.}
  \label{fig:stacking_schematic}
  \includegraphics[width=0.95\textwidth]{../figures/schematic_stacking.png}
\end{figure}

\begin{figure}[H]
  \centering
  \caption{Dose versus realized-dose tag: why the event study can test pre-trends.}
  \label{fig:dose_vs_tag}
  \includegraphics[width=0.95\textwidth]{../figures/schematic_dose_vs_tag.png}
\end{figure}

%==========================================================


\begin{figure}[H]
  \centering
  \caption{Baseline rate level versus cumulative increase, municipal (squares) and school-district (triangles) rates.}
  \label{fig:scatter_level_change}
  \includegraphics[width=0.95\textwidth]{../figures/scatter_level_vs_change_inc_cumulative.png}
\end{figure}

\section{Empirical Strategy}
\label{sec:strategy}
%==========================================================

I estimate two specifications using a continuous stacked
difference-in-differences design. First, I run a pooled-Increase
aggregate (Section~\ref{subsec:pooled-spec}) that collapses all
rate-increase events at each rate-setter into a single only-increase
group. Second, I estimate a binding
decomposition (Section~\ref{subsec:stacked-continuous}) on the three
Increase cells defined in
Section~\ref{sec:theory}, isolating the rate changes that move the
worker's bill from those that do not. Both specifications are
estimated at the municipality-year level. Decrease cells are not estimated: the cohort is
thin ($5$ Dec-bind, $3$ Dec-becomes-non, $11$ Dec-nobind) and the
realized-dose pre-trend rejects parallel trends on every Dec outcome
(Wald $p<0.001$), so the Dec coefficients cannot be interpreted as
causal. It is important to note that a theoretical case for symmetric responses to
rate cuts and rate hikes is important but most likely weak. First, moving costs are
sunk, so the effect of rate cuts may be less salient to incumbent residents
than rate hikes. Pre-trends are tested
separately by a realized-dose Sun--Abraham event study
(Section~\ref{subsec:pretrend-diagnostics}). This section focuses on
estimating equations, identification, and diagnostics.


\subsection{Causal Parameters: ACR, ACRT, and ATT}
\label{subsec:causal-params}
%==========================================================

The headline specification (Equation~\ref{eq:headline}) returns one
coefficient per cell, measured per percentage point. With a binary
treatment that coefficient would be the ATT. With a continuous treatment
the total effect depends on the size of the rate change, so one number
becomes three related objects. Writing the response with a quadratic dose,
$Y_{it} = \alpha_i + \gamma_t + b_1\,\Delta\tau_{it} + b_2\,\Delta\tau_{it}^2 + X_{it}$,
makes them concrete:

\begin{itemize}
  \item $\mathrm{ATT}(d) = b_1 d + b_2 d^2$ is the total effect. A
  municipality that raised its rate by $d$ percentage points ends up
  $\mathrm{ATT}(d)$ log points below a never-changer. 
  \item $\mathrm{ACRT}(d) = b_1 + 2 b_2 d$ is the marginal effect, the
  derivative of $\mathrm{ATT}(d)$. It answers what one more percentage
  point does to a municipality that has already raised its rate by $d$. A
  negative $b_2$ means each additional point matters less than the last.
  \item The ACR is the average of $\mathrm{ACRT}(d)$ over the rate changes
  municipalities actually chose, weighted by how common each dose is. The
  linear specification returns exactly this number
  \citep{bakerDifferenceinDifferencesDesignsPractitioners2025, dechaisemartin2024}.
  It is the headline estimand, and it plays the role for a continuous
  treatment that the ATT plays for a binary one.
\end{itemize}

\noindent The three fit together simply. $\mathrm{ATT}(d)$ accumulates the
marginal effects from $0$ to $d$, $\mathrm{ACRT}(d)$ is its slope, and the
ACR averages that slope over the data.

%==========================================================
\subsection{The Binding Decomposition}
\label{sec:theory}
%==========================================================

The piecewise tax allocation in Equation~(\ref{eq:piecewise}) is the causal
identification backbone of my empirical design. A worker's bill is
$T = \max(\tau^{\text{home}}, \tau^{\text{work,muni}})$. A rate change
only moves the bill when its side of the max binds. The
implied change in the worker's bill from a single rate event is
\[
  \Delta T \;=\; \max(\tau^{\text{home}}_{\text{post}}, \tau^{\text{work,muni}})
              - \max(\tau^{\text{home}}_{\text{pre}},  \tau^{\text{work,muni}}),
\]
which depends jointly on the direction of the rate change and on
which side of the max is binding pre- and post-event. This tax difference from pre and post rate changes can
produce a strong sorting response, a moderate response, or no
response at all depending on the binding transition.

\subsubsection{Six Cells per Rate-Setter}
\label{subsec:cells}

Two definitions structure the decomposition. First, the piecewise
rule $T = \max(\tau^{\text{home}}, \tau^{\text{work,muni}})$ has two
sides: the home side
$\tau^{\text{home}} = \tau^{\text{res,muni}} + \tau^{\text{SD}}$,
set at the worker's residence, and the work side
$\tau^{\text{work,muni}}$, the work rate set by the worker's workplace municipality.
Whichever side is larger is the binding side and determines
the bill. Second, three rate-setters adjust statutory rates
independently across municipality-years: the municipality resident
rate $\tau^{\text{res,muni}}$, the school-district rate
$\tau^{\text{SD}}$, and the municipality work rate
$\tau^{\text{work,muni}}$. The decomposition is built one rate-setter
at a time.

For each rate-setter $\tau$, a cell classifies a municipality
by the direction of its first rate event in $\tau$ (Increase or
Decrease) and then the binding transition of its home stack
across that event: stayed-binding, flipped, or stayed-non-binding.
The cross gives six cells plus a seventh never-changed control pool,
of which only the three Increase cells and the NC control are used in
this paper (Table~\ref{tab:cells}). Each municipality therefore carries three
cell labels (one for each of $\tau^{\text{res,muni}}$,
$\tau^{\text{SD}}$, and $\tau^{\text{work,muni}}$) and the binding
state for each is determined by where its home stack sits relative
to its work rate at the time of that event.

\begin{table}[H]
\centering
\small
\begin{tabular}{lcccl}
\toprule
Cell & Direction & Pre-binding? & Post-binding? & Bill effect \\
\midrule
Inc-bind         & $\uparrow$  & yes & yes & rises by full $\Delta\tau$ \\
Inc-becomes      & $\uparrow$  & no  & yes & rises by part of $\Delta\tau$ \\
Inc-nobind       & $\uparrow$  & no  & no  & no change (placebo) \\
NC               & none & --- & --- & control pool \\
\bottomrule
\end{tabular}
\caption{Binding cells per rate-setter on the increase side. For Inc-becomes the bill rises only by the portion of $\Delta\tau$ above the work rate. Three mirror decrease cells exist by construction but are not estimated in this paper because the decrease cohort is too thin to support causal inference ($5$ Dec-bind, $3$ Dec-becomes-non, $11$ Dec-nobind) and the realized-dose pre-trend rejects parallel trends on every Dec outcome.}
\label{tab:cells}
\end{table}

The cell labels describe the state of the home stack at the
rate event and not necessarily which rate-setter is caused the cell to be that state.
The home stack $\tau^{\text{home}} = \tau^{\text{res,muni}} + \tau^{\text{SD}}$ is
additive in its two home-side components where either component can put
the stack above or below the work floor. A municipality
therefore carries two independent cell labels, one for each
rate-setter (municipality and school district), and the two are not mutually exclusive such as when a muni
classified as Inc-bind on the muni side may be Inc-becomes on the
SD side (or any other combination). The interaction terms in the
estimating equation (Section~\ref{sec:strategy}) separate the
per-pp effect of a muni-rate change from the per-pp effect of an
SD-rate change, each conditional on its own cell label.

What the labels imply for an individual worker depends on where
that worker is employed. Under Act~32 the employer withholds at
$\max(\tau^{\text{home}}, \tau^{\text{work,muni}})$, where
$\tau^{\text{work,muni}}$ is the work rate of the
worker's workplace \citep{paAct32_2008, bagchiEffectsIntroducingWithholding2022}. For a worker who lives and works in
municipality $i$ the workplace rate is exactly $i$'s
own $\tau^{\text{work,muni}}_i$, which is the comparison floor the cell
classification uses. For a resident of $i$ who works in a different
municipality $j$, the relevant floor is $\tau^{\text{work,muni}}_j$,
which the LODES origin--destination data does not let us recover at
the worker level and therefore the label using $i$'s own
work rate is an aggregate proxy that is reasonable on
average across Pennsylvania's local commute shed but not verifiable
individually. This distinction motivates the outcome ordering in
Section~\ref{sec:results}: the count of workers who live and work
in their own municipality, $\log(\text{live-and-work})$, is the
cleanest binding-rule test because the classification matches each
worker's actual workplace and the model directly estimates the response of those who face either $\tau^{\text{home}}$ or $\tau^{\text{work,muni}}$. The other count outcomes are noisier tests: $\log(\text{total jobs})$ and the share of workers living in their own municipality both fold in in-commuters, and $\log(\text{in-commuters})$
does not test the muni resident-rate binding rule at all because
in-commuters' bills do not depend on the resident rate of the
municipality where they work.

The fundamental testable prediction is an
ordered ranking of cell magnitudes where Inc-bind effects exceed
Inc-becomes effects, which in turn exceed Inc-nobind effects, with
the last close to zero and serving as the
within-design placebo. A tax salience argument would state that workers in the Inc-becomes cell
experience a discrete switch from paying zero home-side tax to
paying a positive amount, which should trigger a large behavioral response. The
central empirical comparison of the paper is between Inc-bind and
Inc-nobind where if statutory rates matter mechanically through tax
liability, the residential response should appear in Inc-bind and
not in Inc-nobind. If Inc-nobind moves, the mechanism
interpretation fails and the response is something other than the
binding rule because Inc-nobind events do not change the worker's
bill. This makes the design self-falsifying in a way that a
conventional pooled-Increase specification is not. In the pooled
specification, every increase is a treated event, while here the
institutional rule predicts that only a subset of statutory
increases should generate a response.

%==========================================================
\subsection{Pooled Aggregate Specification}
\label{subsec:pooled-spec}

In this section I report an aggregate specification that pools the
treated Increase cells per rate-setter into a single directional group
$\text{Inc}$ (any rate increase), with never-changers (NC) as the omitted reference. Decrease cells
are excluded because the cohort is
thin ($5$ Dec-bind, $3$ Dec-becomes-non, $11$ Dec-nobind) and the
realized-dose pre-trend rejects parallel trends on every Dec outcome
(Wald $p<0.001$). The pooled
regression answers: what is the average per-percentage-point (pp)
effect of a rate hike on residents, ignoring whether or not the
hike binds? The binding decomposition that follows in
Equation~(\ref{eq:headline}) then explains how the binding mechanism drives the pooled coefficient.

For each unit-time observation $(i, t)$ on the stacked panel, the
per-pp pooled regression is

\begin{equation}
\begin{split}
Y_{\tilde i \tilde t} \;=\; \alpha_{\tilde i} \;+\; \gamma_{\tilde t}
&+ \beta^{M}_{\text{Inc}}\;\Delta\tau^{M}_{it}\,\mathbf{1}\!\left\{d^{M}_{i} = \text{Inc}\right\} \\
&+ \beta^{SD}_{\text{Inc}}\;\Delta\tau^{SD}_{it}\,\mathbf{1}\!\left\{d^{SD}_{i} = \text{Inc}\right\} \\
&+ \mathbf{X}'_{it}\boldsymbol{\phi}
\;+\; \varepsilon_{\tilde i \tilde t}.
\end{split}
\label{eq:pooled}
\end{equation}

The dependent variable $Y_{\tilde i \tilde t}$ is the outcome of
interest. At the municipality level the
outcomes are the logs of live-and-work workers (residents employed
in their own municipality), in-commuters, log total jobs, and the
share of workers living in their own municipality. The unit $i$ is a municipality and $t$ a calendar year. The subscripts $\tilde i = (i, c)$ and
$\tilde t = (t, c)$ are the unit-by-stack and year-by-stack
groupings where each treated unit $i$ appears in exactly one
sub-experiment (its first-event cohort $c$), while each NC unit
appears in every cohort it overlaps. Cohort $c$ is defined as the first year that the rate-setter changed $\tau$. The fixed effects
$\alpha_{\tilde i}$ and $\gamma_{\tilde t}$ are for each unit and cohort pair. The directional grouping
$d^{s}_{i} \in \{\text{NC}, \text{Inc}\}$ classifies
unit $i$'s rate trajectory on rate-setter $s \in \{M, SD\}$, with
NC the omitted reference as they are the control group. The
dose $\Delta\tau^{s}_{it} = \tau^{s}_{it} - \tau^{s}_{i, c-1}$ is
the cumulative change in the rate from the pre-event baseline
($c-1$, one year before the first event) to year $t$, measured in
percentage points; it is zero in pre-event years and for NC units,
and accumulates with subsequent rate changes within $i$. The control vector
$\mathbf{X}_{it}$ holds the six municipality fiscal share controls
(police-services, streets, debt-service, EIT-revenue,
real-estate-tax, intergovernmental-revenue).
$\boldsymbol{\phi}$ is the
coefficient vector on these controls and $\varepsilon_{\tilde i \tilde t}$
is the error term.

\subsection{Binding Decomposition: Stacked Continuous DiD}
\label{subsec:stacked-continuous}

The binding decomposition is the main estimator of the paper. It
replaces the two-level direction grouping $d^s_i$ from
Equation~\ref{eq:pooled} with a multilevel cell
$\text{cell}^s_i$ that entails both the direction of unit $i$'s
first event and whether the home stack was binding pre- and
post-event (Section~\ref{subsec:cells}, Table~\ref{tab:cells}).
This splits the increase side into Inc-bind, Inc-becomes, and
Inc-nobind, while NC remains the omitted reference group. Decrease cells are
not estimated for the reasons discussed in
\ref{sec:strategy}. The stacked
panel, unit-by-stack groupings, dose $\Delta\tau^s_{it}$, controls
$\mathbf{X}_{it}$, fixed effects, and policy-jurisdiction
clustering are identical to Equation~\ref{eq:pooled}. The
sub-experiment for each cohort $c$ runs on the window
$t \in [c-K, c+K]$ with $K=5$, which gives the realized-dose ES
diagnostic four pre-event leads (Section~\ref{subsec:pretrend-diagnostics})
and covers the bulk of the treated units with a full $5$-pre,
$5$-post window.

For each rate-setter $s$ and outcome $Y$ the stacked-continuous
regression is

\begin{equation}
\begin{split}
Y_{\tilde{i}\tilde{t}} \;=\; \alpha_{\tilde{i}} \;+\; \gamma_{\tilde{t}}
&+ \sum_{k \in \mathcal{C}_{-\text{NC}}} \beta^{\text{muni}}_k \;\Delta\tau^{\text{muni}}_{it} \cdot \mathbf{1}\{\text{cell}^{\text{muni}}_i = k\} \\
&+ \sum_{k \in \mathcal{C}_{-\text{NC}}} \beta^{\text{SD}}_k  \;\Delta\tau^{\text{SD}}_{it}   \cdot \mathbf{1}\{\text{cell}^{\text{SD}}_i = k\}   \\
&+ \mathbf{X}'_{it}\boldsymbol{\phi}
\;+\; \varepsilon_{\tilde{i}\tilde{t}}.
\end{split}
\label{eq:headline}
\end{equation}

The set
$\mathcal{C}_{-\text{NC}} = \{\text{Inc-bind, Inc-becomes, Inc-nobind}\}$
collects the three estimated treated cells; $\text{cell}^s_i$ is the cell label
on rate-setter $s$ and the unit level being modeled $i$. The regression returns three $\beta^s_k$
coefficients per rate-setter. Because the dose $\Delta\tau^s_{it}$ enters
linearly, each $\beta^s_k$ is an average causal response (ACR): the
dose-weighted average per-percentage-point effect of rate $s$ on $Y$ within
cell $k$, not the marginal effect at any single dose
\citep{bakerDifferenceinDifferencesDesignsPractitioners2025,dechaisemartin2024}.
Section~\ref{subsec:rob-acr} shows this ACR is recovered with non-negative
weights and is not distorted by dose curvature. The
$\beta^{\text{muni}}_{\text{Inc-nobind}}$ coefficient is the
within-design placebo meaning that the rate hike
does not move the worker's bill in this cell, so the coefficient
should be statistically indistinguishable from zero. All other
elements of the regression are inherited from
Equation~\ref{eq:pooled}.

The construction uses the negative-weighting fix from
\citet{bakerDifferenceinDifferencesDesignsPractitioners2025}:
treated units enter exactly one stack (their first-event cohort),
NC units enter every stack, so no treated unit serves as control
to a later-treated unit. The cell decomposition is a
novel addition because it generates a ladder of predictable effects including the within-design placebo
(Inc-nobind), which the pooled stacked DiD
classification cannot produce.

The identifying variation in this design comes from the cross-unit cumulative rate changes after each unit's first
event. Figure~\ref{fig:dose_trajectory} and Figure~\ref{fig:dose_trajectory_inc_becomes}
show the average change post event by rate-setter. Each line plots the cross-muni mean and
inter-quartile range of cumulative $\Delta\tau$ at each event-time. Figure~\ref{fig:dose_trajectory}
traces the mean cumulative resident-rate change of Inc-bind units relative to their pre-event baseline. The
pre-event window is mechanically flat at zero by construction (dose
is measured relative to the unit's own $c-1$ rate) and the post-event
window shows the realized cumulative $\Delta\tau^{\text{muni}}$
rising to roughly $0.24$ percentage points by $k = +5$. While the school district rate sits far above the municipal rate at roughly $0.60$ percentage points.


\begin{figure}[H]
  \centering
  \caption{Cumulative Inc-bind dose by event-time: municipality cohort vs.\ school-district cohort.}
  \label{fig:dose_trajectory}
  \includegraphics[width=0.85\textwidth]{../figures/support_dose_trajectory_muni_vs_sd.png}
  \begin{minipage}{0.85\textwidth}
    \footnotesize \textit{Notes:} Mean and inter-quartile range of
    cumulative dose by relative event-time $k$. The municipality
    cohort (red) reaches a mean of $0.24$ pp at $k=+5$; the
    school-district cohort (blue) reaches $0.65$ pp. Pre-event
    values are zero by construction.
  \end{minipage}
\end{figure}

Figure~\ref{fig:dose_trajectory_inc_becomes}
traces the mean cumulative resident-rate change of Inc-becomes binding on units relative to their pre-event baseline. The post-event
window shows a similar effect as figure~\ref{fig:dose_trajectory} where the school district rates sit above municipality rates at $0.60$ percentage points but the municipality rate above $0.40$ which is a much larger increase than the inc-binds grouping thus shrinking the gap between the average school district and average municipality rate. Figure~\ref{fig:dose_trajectory_inc_becomes} indicates that for an the rate to become binding the jump may primarily due to municipality rate adjustments.

\begin{figure}[H]
  \centering
  \caption{Cumulative Inc-becomes dose by event-time: municipality cohort vs.\ school-district cohort.}
  \label{fig:dose_trajectory_inc_becomes}
  \includegraphics[width=0.85\textwidth]{../figures/support_dose_trajectory_inc_becomes.png}
  \begin{minipage}{0.85\textwidth}
    \footnotesize \textit{Notes:} As in Figure~\ref{fig:dose_trajectory}
    but for the Inc-becomes cell, where the home stack flips from
    non-binding to binding at the first event.
  \end{minipage}
\end{figure}

\subsection{Realized-Dose Event-Study Diagnostic}
\label{subsec:pretrend-diagnostics}
\label{subsec:realized-dose-es}

Equations~\ref{eq:pooled} and~\ref{eq:headline}  cannot test for pre-treatment leads
directly because $\Delta\tau^s_{it}$ is mechanically zero for all
$t < c$, so any pre-event coefficient on the dose
is also mechanically zero. I therefore run a separate
realized-dose Sun--Abraham event study on the same stacked panel
to test parallel trends, the condition recommended
for continuous-treatment DiD designs by
\citet{callaway2021} and
\citet{dechaisemartin2024} states that units with
larger eventual rate changes should not be on differentially sloped
pre-event trajectories relative to NC controls. For each cell on
rate-setter $s$ I restrict the panel to units in (NC, cell) and
run

\begin{equation}
Y_{it} = \alpha_{\tilde i} + \gamma_{\tilde t}
+ \sum_{k \neq -1} \beta^M_k\, \mathbf{1}\{\ell_{it} = k\}\,
\widetilde{D}_i
+ \mathbf{X}'_{it}\boldsymbol{\phi} + \varepsilon_{it}.
\label{eq:realized-dose-es}
\end{equation}

Two variables are new relative to Equation~\ref{eq:headline}. The
event-time index $\ell_{it} = t - c^M_i$ counts the years between
calendar year $t$ and unit $i$'s first muni-rate event year $c^M_i$,
with $\ell = -1$ (one year before the event) as the omitted
reference; in the sum below $k$ denotes a value of $\ell$, i.e.\ an
event-time, in contrast to Equation~\ref{eq:headline}, where $k$
indexed the binding cells. The realized-dose tag
$\widetilde{D}_i = \tau^M_{i, c+5} - \tau^M_{i, c-1}$ is a single number per unit where $i$'s eventual cumulative rate
change is the observed rate five years after its first event minus the
rate the year before and held fixed across all event-time years rather
than allowed to grow from zero as $\Delta\tau^M_{it}$ does in
Equation~\ref{eq:headline}. It enters the regression only through the
interaction $\mathbf{1}\{\ell_{it}=k\}\,\widetilde{D}_i$, which equals
$\widetilde{D}_i$ in the single year unit $i$ is at event-time $k$ and
zero in all its other years because that varies within unit, the
unit fixed effect (which removes only within-unit constants) does not
absorb it. Because the tag is non-zero in every event year and varies across units by eventual treatment
intensity, each pre-event $\beta^M_k$ tests whether units that
eventually take larger rate changes were already on differently sloped
trajectories. The time-varying $\Delta\tau^M_{it}$ cannot do this: it
is mechanically zero before the event, so its pre-event coefficients
are mechanically zero. The outcome $Y_{it}$, the unit-by-stack and year-by-stack
fixed effects, the control vector $\mathbf{X}_{it}$, the units of
observation, and the policy-jurisdiction clustering are the same
as in Equation~\ref{eq:headline}.

The pre-event coefficients $\beta^M_k$ for $k \leq -2$ are the
parallel-trends test: if units that eventually take large rate
changes are on differentially sloped trajectories before the rate
change, the $\beta^M_k$ at pre-event leads will be non-zero. Under
strong parallel trends they are zero. The post-event coefficients
$\beta^M_k$ for $k \geq 0$ recover the per-pp dose-response at
each event-time and average to the corresponding cell coefficient
$\beta^s_{\text{cell}}$ from Equation~\ref{eq:headline} when the
unit's rate path is a single step at the event year. I report the
joint Wald $p$-value as a single pre-trend
statistic for each cell-outcome combination.

Section~\ref{sec:results} reports the headline binding
decomposition at the municipality and the pooled-Increase
benchmark. Additional diagnostics (pre/post Act~32 split, Act~47
exclusion) are in Section~\ref{sec:robustness}.

%==========================================================
\section{Results}
\label{sec:results}
%==========================================================

I report results in two layers. First,
Section~\ref{subsec:results-pooled} presents the pooled Inc
benchmark, which approximates what a conventional local-tax design
would recover by averaging across all rate-increase events without
distinguishing whether they bind. Second,
Section~\ref{subsec:results-decomp} reports the binding
decomposition at the municipality level. The central empirical
comparison is between the Inc-bind cell (where the statutory rate
change moves the worker's bill) and the Inc-nobind cell (where it
does not), with Inc-nobind serving as the within-design placebo.
Dec cells are not estimated: the cohort is thin
(5 Dec-bind, 3 Dec-becomes-non, 11 Dec-nobind) and the
realized-dose pre-trend rejects parallel trends on every Dec
outcome (Wald $p<0.001$), so the Dec coefficients cannot be
interpreted as causal. The pre-/post-Act-32 enforcement-break split
appears in the Robustness section.

\subsection{Pooled Results}
\label{subsec:results-pooled}

I first report the pooled-Increase specification from
Equation~(\ref{eq:pooled}), which collapses each rate-setter to a
single direction tag for rate increases. The pooled coefficients
are useful because they approximate the estimand in a conventional
local-tax design where the average per-pp dose-response of a
municipality-rate change and an SD-rate change without
distinguishing whether the rate change moves the worker's bill.

Table~\ref{tab:pooled_per_pp_cousub} reports the per-pp pooled
coefficients at the municipality level. The table reports the two
coefficients of interest per outcome: the municipality resident-rate
Inc effect $\beta^M_{\text{Inc}}$ and the school-district-rate Inc
analogue $\beta^{SD}_{\text{Inc}}$.

\input{../tables/pooled_per_pp_cousub.tex}

% [Results prose removed 2026-07-02 pending property-controls re-run;
%  regenerate numbers from data/q3_unified CSVs before rewriting.]

\begin{figure}[H]
  \centering
  \caption{Dose-response: ATT$(d)$ and ACRT$(d)$ with the linear ACR (pooled Increase).}
  \label{fig:spline_pooled}
  \includegraphics[width=0.95\textwidth]{../figures/acr_acrt_att_spline_pooled.png}
\end{figure}

\subsection{Decomposition Results}
\label{subsec:results-decomp}
\label{subsec:results-headline}

The decomposition is estimated at the municipality level. Both
rate-setter treatments enter the same joint regression: each
muni-year is assigned its own resident rate $\tau^M$ at local
level and $\max(\tau^{SD})$ across the school districts the muni
reports to. Table~\ref{tab:headline_muni_combined}
reports both sets of cell coefficients from the same joint
regression in a single panel: the muni resident-rate cells (top
panel) are the headline causal claim, and the SD resident-rate
cells (bottom panel) are reported alongside as co-treatment from
the identical specification.

\input{../tables/headline_muni_combined.tex}

% [Results prose removed 2026-07-02 pending property-controls re-run.
%  Structure to restore: typical-dose translation (Inc-bind), share outcome,
%  Inc-becomes read, Inc-nobind placebo, SD-side cells.]

\begin{figure}[H]
  \centering
  \caption{Dose-response: ATT$(d)$ and ACRT$(d)$ with the linear ACR (Inc-bind cell).}
  \label{fig:spline_incbind}
  \includegraphics[width=0.95\textwidth]{../figures/acr_acrt_att_spline_incbind.png}
\end{figure}

\begin{figure}[H]
  \centering
  \caption{Dose-response: ATT$(d)$ and ACRT$(d)$ with the linear ACR (Inc-becomes cell).}
  \label{fig:spline_incbecomes}
  \includegraphics[width=0.95\textwidth]{../figures/acr_acrt_att_spline_incbecomes.png}
\end{figure}

\subsection{Event-Study Results}
\label{subsec:results-es}

% ---- ES diagnostic tables moved out of main text (full lead/lag tables in the
% ---- Event-Study appendix; the overlay figure below is the main-text exhibit):
% \subsubsection{Realized-dose ES benchmark diagnostic}
% \label{subsubsec:pooled-es}
%
% Table~\ref{tab:es_pooled_cousub} reports the ATT summary of the
% muni-level pooled realized-dose event study, with the
% cluster-robust Wald pre-trend $p$-value for the joint null
% $\beta_k = 0$ at $k \leq -2$.
% Figure~\ref{fig:pooled_es_muni_only} plots the event-time path.
%
% \input{../tables/es_pooled_cousub_summary.tex}
%
% % [Results prose removed 2026-07-02 pending property-controls re-run.]
%
% \begin{figure}[H]
%   \centering
%   \caption{Municipality-level pooled realized-dose event study.}
%   \label{fig:pooled_es_muni_only}
%   \includegraphics[width=0.95\textwidth]{../figures/pooled_es_muni_only.png}
%   \begin{minipage}{0.95\textwidth}
%     \footnotesize \textit{Notes:} Per-pp realized-dose coefficient at each event-time for the pooled Inc cohort. Ribbon is the 95\% CI; $k=-1$ omitted as the reference period.
%   \end{minipage}
% \end{figure}
%
% \subsubsection{Realized-dose ES decomposition diagnostics}
% \label{subsubsec:decomp-es}
%
% The realized-dose event study disaggregates the pooled pre-trend
% by cell. Table~\ref{tab:es_realized_cousub} reports each cell's
% post-event ATT, pre-event lead, and the joint pre-trend Wald $p$
% across the four muni outcomes.
% % [Results prose removed 2026-07-02 pending property-controls re-run.]
%
% \input{../tables/es_realized_cousub.tex}
%
% \begin{figure}[H]
%   \centering
%   \caption{Municipality-level realized-dose event study on log live-and-work, by cell.}
%   \label{fig:es_realized_muni_live_work}
%   \includegraphics[width=0.95\textwidth]{../figures/es_realized_log_live_and_work.png}
%   \begin{minipage}{0.95\textwidth}
%     \footnotesize \textit{Notes:} Per-pp coefficient at each event-time, faceted by cell.
%   \end{minipage}
% \end{figure}
%
% %==========================================================
%
\begin{figure}[H]
  \centering
  \caption{Realized-dose event study by group (pooled / Inc-bind / Inc-becomes), level-at-mean-dose scale.}
  \label{fig:feols_es_bycell}
  \includegraphics[width=0.95\textwidth]{../figures/feols_es_bycell.png}
\end{figure}

% ---- EIT Revenue Response + Summary subsections commented out
% ---- per outline restructure 2026-07-06 (models still run in stage 02):
% \subsection{EIT Revenue Response}
% \label{subsec:results-revenue}
% 
% Sorting erodes the tax base, so the revenue raised by a rate increase is
% less than the mechanical projection.
% % [Results prose removed 2026-07-02 pending property-controls re-run.]
% 
% %==========================================================
% \subsection{Summary}
% \label{subsec:results-summary}
% 
% % [Results prose removed 2026-07-02 pending property-controls re-run.]
% 
\section{Robustness}
\label{sec:robustness}
%==========================================================

This section summarizes the robustness checks behind the headline
muni-level dose-response. The realized-dose event-study pre-trend
diagnostic was discussed alongside the headline in
Section~\ref{subsubsec:decomp-es}. Full tables behind the
robustness summaries are in Appendix~\ref{app:robustness}.

\subsection{de Chaisemartin and D'Haultfœuille Estimator}
\label{subsec:rob-acr}

Each headline $\beta^M_k$ is an ACR, and a linear-in-dose stacked
specification recovers it cleanly only if the implicit weights are
non-negative and the dose-response is not strongly curved
\citep{dechaisemartin2024}. I test both concerns by re-estimating the
muni-side increase response on the same sample with the non-negative-weight
continuous-treatment estimator of
\citet{chaisemartinUsingDid_multiplegt_dynEstimate2025}.
% [Results prose removed 2026-07-02 pending property-controls re-run.]

\input{../tables/dcdh_results.tex}

\noindent Reading Table~\ref{tab:dcdh_results}, per group:
\begin{itemize}
  \item The first two rows are the estimated effect at $k=+5$ and its
  standard error: the total change in the outcome five years after the
  rate change, at the doses municipalities actually took. These are level
  effects, so compare them to the linear ACR times the group's mean dose.
  \item Dep.\ mean is the average outcome level in that group's
  estimation sample, and it converts effects into head counts. A log
  live-and-work mean of $3.7$ is roughly $41$ workers, so an effect of
  $-0.07$ log points (about a $7\%$ decline) is roughly three fewer
  live-and-work workers in the typical municipality.
  \item Avg total effect per pp is the estimator's dose-normalized
  average. It sits on the same per-pp scale as the headline ACR, so the
  two numbers can be compared directly.
  \item Joint placebo $p$ tests for pre-trends. Values above conventional
  levels mean treated and control municipalities were on similar paths
  before the change, so larger is better.
  \item $N$ and Switchers count the observations and treated
  municipalities identifying the $k=+5$ effect.
\end{itemize}

\begin{figure}[H]
  \centering
  \caption{dCDH event study by group (non-negative weights): pooled, Inc-bind, Inc-becomes.}
  \label{fig:dcdh_bycell}
  \includegraphics[width=0.95\textwidth]{../figures/dcdh_bycell_es.png}
\end{figure}



\subsection{Pre/Post Act 32 Decomposition}
\label{subsec:rob-act32}

Splitting the headline cell-$\times$-dose specification at the Act-32
enforcement break (cohort first event year $<2012$ vs $\geq 2012$)
tests whether the binding-cell effect strengthens under consolidated
employer withholding
(\hyperref[app:rob-act32]{see Appendix~\ref{app:rob-act32}}).

% \input{../tables/act32_cousub.tex}  % concise version
\input{../tables/act32_full_pre.tex}
\input{../tables/act32_full_post.tex}
% [Results prose removed 2026-07-02 pending property-controls re-run.]

\subsection{Act 47 Distressed-Municipality Exclusion}
\label{subsec:rob-act47}

The Inc-bind cohort includes Act~47 distressed cities
\citep{paAct47_1987}. Act~47 is the legal designation that
lets a municipality exceed the $0.5\%$ statutory municipality-EIT
cap meaning there may be a selection confounder. An
Act~47 designation may come from a period of population
decline, (rate hike occurs after the municipality is
already losing residents), debt
restructuring, and service cuts. Only
$2$ of the $24$ Inc-bind municipalities match the DCED Act~47
designation list (Carbondale and Hazleton). The other historically
Act-47-designated PA cities (e.g.\ Braddock, Coatesville, Johnstown)
do not appear in the Inc-bind cell here because their rate trajectories are non-monotone (did not move in one direction) and were removed (Section~\ref{subsec:stacked-continuous}).
Table~\ref{tab:act47_robustness} re-estimates the
municipality-level Inc-bind coefficient on the remaining $22$
non-Act-47 Inc-bind municipalities.

% \input{../tables/act47_robustness.tex}  % concise 2-col version
\input{../tables/act47_full_results.tex}

% [Results prose removed 2026-07-02 pending property-controls re-run.]

\subsection{Pre-Trend Battery}
\label{subsec:rob-battery}

The joint Wald test over-rejects on cohorts of 14--24 treated
municipalities, so Table~\ref{tab:pretrend_battery} reports four
pre-trend tests per cell: the joint Wald on the realized-dose leads,
the count of individually significant leads, a one-degree-of-freedom
linear pre-slope, and the HonestDiD breakdown bound.
% [Prose TO WRITE after reviewing the table.]

\input{../tables/pretrend_battery.tex}

\subsection{Dose-Permutation Randomization Inference}
\label{subsec:rob-permutation}

Strong parallel trends is a statement about dose selection. This test
shuffles the realized dose magnitudes across treated municipalities,
holding each municipality's event timing and binding cell fixed, and
re-estimates the decomposition on each shuffle.
% [Prose TO WRITE after reviewing the table and figure.]

\IfFileExists{../tables/perm_inference.tex}{\input{../tables/perm_inference.tex}}{\textit{[Permutation run in progress.]}}

\IfFileExists{../figures/perm_inference_hist.png}{%
\begin{figure}[H]
  \centering
  \caption{Dose-permutation distributions with the true estimates.}
  \label{fig:perm_inference}
  \includegraphics[width=0.95\textwidth]{../figures/perm_inference_hist.png}
\end{figure}}{}

\subsection{Fake-Timing Placebo}
\label{subsec:rob-faketiming}

Each treated municipality is assigned a placebo event four years before
its real one, and only data from before the real event enter the
estimation. Nothing happened at the placebo dates, so the coefficients
should be null.
% [Prose TO WRITE after reviewing the table.]

\IfFileExists{../tables/fake_timing.tex}{\input{../tables/fake_timing.tex}}{\textit{[Fake-timing run in progress.]}}

% \section*{Acknowledgments}
% TODO: fill in before submission

%==========================================================
% BIBLIOGRAPHY
%==========================================================
%==========================================================
\section{Discussion}
\label{sec:discussion}
% [TO WRITE -- lit comparison, follows-the-worker vs capitalizes mechanism,
%  binding cell isolates the bill, policy, limits]

\section{Conclusion}
\label{sec:conclusion}
% [TO WRITE -- restate question + headline; no new results]

\newpage
\printbibliography


%==========================================================
%==========================================================
\appendix

\section{Spline Specification for the Dose-Response Figures}
\label{app:spline-spec}

\begin{figure}[H]
  \centering
  \caption{Treated doses and knot placement, by estimation group.}
  \label{fig:knot_placement}
  \includegraphics[width=0.95\textwidth]{../figures/spline_knot_placement.png}
\end{figure}
% [Spec equations from outline to add at write time: pooled + per-cell
%  muni and SD splines; knots at dose terciles; ATT(d), ACRT(d) formulas;
%  BIC selects linear; spline is for visuals.]

\section{Curvature Diagnostics}
\label{app:curvature}

\begin{figure}[htbp]
\centering
\includegraphics[width=0.92\textwidth]{../figures/curvature_att_curve.png}
\caption{Quadratic ATT$(d)$ dose-response by outcome (pooled Increase, headline
sample). Each panel plots $\widehat{\mathrm{ATT}}(d)=b_1 d + b_2 d^2$, the total
effect on $Y$ of a cumulative dose-$d$ rate increase relative to never-changers,
with a delta-method 95\% confidence band. The $d^2$ term is significant for the
two residence outcomes (panel headers report its $p$-value), producing a concave,
diminishing-marginal dose-response; it is insignificant for the job-flow outcomes,
where the response is effectively linear.}
\label{fig:curvature_att}
\end{figure}

\section{Full Results with Control Coefficients}
\label{app:full-results}
%==========================================================

This appendix reports the muni-level binding-decomposition
regression with the fiscal-control coefficients and diagnostics.
Table~\ref{tab:app_full_muni_cell_results} reports the muni-rate
and SD-rate cells (two panels) with the shared fiscal controls from
the same joint regression.

\input{../tables/app_full_muni_cell_results.tex}

%==========================================================
\section{Realized-Dose Event Studies, Full Set}
\label{app:es-full}

\input{../tables/es_full_pooled.tex}
\input{../tables/es_full_incbind.tex}
\input{../tables/es_full_incbecomes.tex}
\input{../tables/es_full_incnobind.tex}

\input{../tables/es_pooled_cousub_summary.tex}
\input{../tables/es_realized_cousub.tex}

\label{app:es-realized}
%==========================================================

Per-cell realized-dose Sun--Abraham event studies at the muni
level for the remaining residence-side outcomes (log live-and-work
is in the body, Section~\ref{subsubsec:decomp-es}).

\begin{figure}[H]
  \centering
  \caption{Muni-level realized-dose ES on $\log(\text{in-commuters})$, by cell.}
  \label{fig:app_es_realized_muni_incomm}
  \includegraphics[width=0.95\textwidth]{../figures/es_realized_log_in_commuters.png}
\end{figure}

\begin{figure}[H]
  \centering
  \caption{Muni-level realized-dose ES on $\log(\text{total jobs})$, by cell.}
  \label{fig:app_es_realized_muni_jobs}
  \includegraphics[width=0.95\textwidth]{../figures/es_realized_log_total_jobs.png}
\end{figure}

\begin{figure}[H]
  \centering
  \caption{Muni-level realized-dose ES on the share of workers
           living in their own municipality, by cell.}
  \label{fig:app_es_realized_muni_share_in}
  \includegraphics[width=0.95\textwidth]{../figures/es_realized_share_live_in.png}
\end{figure}

%==========================================================
\section{Robustness: Full Tables}
\label{app:robustness}
%==========================================================

Full tables behind the robustness summaries in
Section~\ref{sec:robustness}.

\subsection{Pre/post Act 32 split}
\label{app:rob-act32}

Headline cell $\times$ dose specification re-estimated on pre-Act-32
($c<2012$) and post-Act-32 ($c\geq 2012$) cohort subsamples at the
municipality level.

\input{../tables/act32_cousub.tex}  % concise pre/post summary (full versions in Robustness)

%==========================================================
\end{document}
